\documentclass[11pt]{article}

\usepackage[margin=1in]{geometry}
\usepackage[T1]{fontenc}
\usepackage{lmodern}
\usepackage{amsmath,amssymb,mathtools}
\usepackage{booktabs,tabularx,longtable,array}
\usepackage[numbers,sort&compress]{natbib}
\usepackage{microtype}
\usepackage{xcolor}
\usepackage{hyperref}
\usepackage{lineno}
\usepackage{setspace}

\definecolor{icpnavy}{HTML}{183B56}
\definecolor{icpteal}{HTML}{2B6F77}

\hypersetup{
  colorlinks=true,
  linkcolor=icpnavy,
  citecolor=icpteal,
  urlcolor=icpteal,
  pdftitle={The Constrained Future: A Testable Information-Constraint Principle for Persistent Organization},
  pdfauthor={Douglas Clark},
  pdfsubject={BioSystems hypothesis paper submission},
  pdfkeywords={information constraint, adjacent possible, causal intervention, persistent organization, reachable futures, self-organization, relative entropy}
}

\setlength{\tabcolsep}{5pt}
\renewcommand{\arraystretch}{1.14}
\emergencystretch=2em
\onehalfspacing

\newcommand{\ICP}{\textsc{ICP}}
\newcommand{\IOTA}{\textsc{IOTA}}
\newcommand{\Efamily}{\mathcal{E}_{\tau}}
\newcommand{\Rset}{\mathcal{R}}
\newcommand{\status}[1]{\textsf{#1}}

\title{\textbf{The Constrained Future}\\
\large A Testable Information-Constraint Principle for Persistent Organization}
\author{Douglas Clark\\
\small Independent Researcher\\
\small Corresponding author: \href{mailto:clarkdd@gmail.com}{clarkdd@gmail.com}}
\date{}

\begin{document}
\maketitle

\noindent\textbf{Article type:} Hypothesis Paper

\begin{abstract}
Information in living and artificial systems is often quantified as uncertainty,
yet inherited organization also changes which futures remain accessible. We
formulate the Information-Constraint Principle (ICP): physically effective
information is an interaction-formed, embodied constraint on an explicitly
indexed family of reachable or probable future responses, persisting only when
its causal role is re-instantiated in a successor carrier. A critical synthesis
of control and viability theory, computational mechanics, process tensors,
counterfactual semantic information, biological organization, the adjacent
possible, and emergent quasiparticles shows that most future-option language is
established prior art. We therefore specify a persistent-carrier test requiring
formation, sufficiency, intervention, replacement, and representation
robustness. The adjacent possible is formalized as the novelty-bearing type
projection of a policy- and horizon-indexed reachable set. For nested finite
uniform supports, the signed deformation
$\chi(R_1\!:\!R_0)=\log[\mu(R_0)/\mu(R_1)]$ equals forward Kullback--Leibler
divergence under pruning and negative reverse divergence under expansion;
enabled, lost, and retained identities must accompany the scalar for non-nested
supports. Bounded IOTA interventions demonstrate carrier replacement and finite
opportunity windows but not open-ended state-space creation. ICP is therefore a
falsifiable hypothesis rather than a new law: its decisive test is whether
carrier interventions predict held-out future-law changes beyond complete
native dynamics, including a preregistered crossover from immediate microstate
pruning to longer-horizon role expansion. A Gravity-from-Entropy
log-determinant mapping is retained only as an untested outlook.
\end{abstract}

\noindent\textbf{Keywords:} information constraint; adjacent possible; causal
intervention; persistent organization; reachable futures; self-organization;
relative entropy

\linenumbers

\section{The persistence problem}

The motivating question is elementary to state and difficult to answer without
slipping between meanings of information: how can the consequence of an
interaction persist after the object or arrangement that first carried it has
changed or disappeared? A footprint can outlast a foot; a measurement record
can be copied into an environment; angular momentum can be transferred among
spin, orbital, and field degrees of freedom; a biological response can survive
component turnover. None of these examples requires information to float free
of physical realization. They instead distinguish persistence beyond an
\emph{original} carrier from existence without \emph{any} carrier.

IOTA takes that distinction as a research problem rather than a slogan. Its
organizing hypothesis is that information is the interaction-caused change in
the reachable. In condition form, the information present at time $t$ is the
part of realized interaction history that remains causally active in present
local state and relational topology. In change form, an interaction is
informational when it changes which future trajectories or responses are
reachable, or changes their probabilities. The proposal is intentionally
stronger than Shannon information as uncertainty and narrower than the claim
that every physical difference is information.

This paper gives the synthesis two jobs. The first is diagnostic: does this
program have genuine scientific legs when compared with the best existing
science, including theories that may already contain its central object? The
second is constructive: if it does, what sequence of results would move it from
a bounded model of persistent constraint toward an account of familiar physics?
The review is therefore allowed to return a negative verdict. A literature
collection selected only because it sounds information-friendly would defeat
both purposes.

\subsection{IOTA as substrate hypothesis, not a frozen cosmic lattice}

The IOTA substrate hypothesis treats informational events or distinctions,
relations, and local rules of lawful succession as the irreducible explanatory
layer. The concrete lattice in the present simulation is a candidate
realization, not a claim that its dimensionality, neighborhood, scheduler, or
constants are the uniquely correct microscopic architecture of nature. This
distinction is essential. A literal fixed lattice inherits immediate burdens:
preferred structure, continuum and symmetry recovery, dependence on update
order, and the possibility that every reported coherence is a special feature
of one engineered rule.

At the theoretical level, write the retained present condition as
\begin{equation}
 H_{\le t}\longrightarrow M_t=(X_t,G_t),
 \label{eq:history-record}
\end{equation}
where $H_{\le t}$ is realized interaction history, $X_t$ the current local
state, and $G_t$ the active relational organization. Fixed law, topology,
scheduler, boundary, constants, and exogenous inputs must be declared rather
than hidden in the word ``system.'' Two histories are operationally equivalent
at $t$ only when no allowed future intervention distinguishes the retained
states at the chosen scale and precision.

The ontology claim is a postulate under investigation: familiar physical
objects and quantities are proposed to be stable effective organizations of
this underlying structure. It is not inferred from the mere use of a computer
model, nor from the truism that physical theories can be encoded as data.

\subsection{ICP and the emergence ladder}

ICP is the proposed principle connecting the substrate hypothesis to dynamics:

\begin{quote}
\textbf{Physically effective information is interaction-formed, embodied
constraint on an explicitly indexed family of reachable or probable futures;
it persists when that constraint role is lawfully re-instantiated beyond its
original carrier.}
\end{quote}

The proposal has a definite explanatory order:
\begin{enumerate}
  \item local distinctions and relations make some successor events lawful and
  others unavailable;
  \item interaction changes that successor structure;
  \item some organizations preserve and transmit selected constraints;
  \item carrier succession allows a relational or causal role to outlast its
  original realization;
  \item persistent coherences may acquire stable collective invariants;
  \item suitable coherences may realize quasiparticle and field roles; and
  \item mass, energy, geometry, clocks, and macroscopic physics may appear as
  effective properties of those organizations.
\end{enumerate}

Only the beginning of this ladder has project evidence. The current quantum
Dirac construction starts higher up: it assumes quantum amplitudes,
composition, a Born readout, a clock, and a mass schedule, then tests whether
specified local updates recover a bounded known role
\citep{Clark2026DiracRole}. It is evidence that a local construction can realize
that role under supplied quantum structure, not yet a derivation of quantum
quasiparticles from the base IOTA lattice. The gap is not a footnote. It is the
central research frontier.

\subsection{Why the synthesis must be adversarial}

The phrase ``constrains the future'' is dangerously easy to satisfy. Given a
complete physical state and law, virtually every state difference can alter a
future trajectory. Control theory already studies reachable sets; statistical
models already condition future distributions; causal models already compare
interventions; and semantic-information work already erases correlations and
measures effects on later viability. If ICP merely renames these objects, it may
be useful language but is not a new physical principle.

The burden is therefore not to find fields in which constraints matter. It is
to determine whether the combined requirements of historical formation,
physical embodiment, sufficiency for a declared intervention-response family,
causal carrier manipulation, succession across original-carrier replacement,
and representation robustness produce a new invariant, compression,
prohibition, or prediction.

\section{Method and evidence discipline}

This article is a \emph{critical integrative synthesis and operational
proposal}. It is not a systematic review or a claim of exhaustive retrieval.
The project registry records the query strata, primary-source checks, source
notes, competing explanations, falsification criteria, and a claim--evidence
matrix. Searches were expanded adversarially whenever a field offered a likely
redundancy account. The final bounded closure covered classical
reachability/viability, causal and semantic information, feedback information
thermodynamics, positive and critical Quantum Darwinism, quantum reference
frames, biological organization and carrier interventions, and emergent
quasiparticles. The bounded cross-domain search cutoff was 20 July 2026. A
targeted 21 July amendment added two primary Gravity from Entropy sources and
the formal IOTA--GfE bridge discussed below. A second question-triggered
amendment on 21 July reviewed the adjacent possible, correlated-novelty models,
and the Theory of the Adjacent Possible (TAP). Neither amendment reopened the
other search strata or converted the review into an exhaustive search.

Claims are assigned to five evidence strata:
\begin{enumerate}
  \item \textbf{Established field result}: a theorem, experiment, or formal
  construction within its native assumptions.
  \item \textbf{Project constructive evidence}: a frozen or prospectively
  registered construction within a declared model.
  \item \textbf{Project empirical/computational evidence}: measurements on an
  engineered or observed data-generating system, without reclassifying that
  system as fundamental physics.
  \item \textbf{Synthesis inference}: a cross-field structural mapping that may
  be supported even when no source states it.
  \item \textbf{IOTA/ICP hypothesis}: a proposed universality, emergence, or
  ontology claim.
\end{enumerate}

Publication and peer review are relevant quality signals, not truth labels.
The IOTA and Dirac manuscripts are treated as project evidence; the PSA
semantic-field manuscript is under review. None is counted as independent
replication. Positive, negative, stopped, censored, and withdrawn project
results are retained because a viability audit cannot be built only from
successful gates.

\subsection{Indexed future-law families}

For a candidate retained state or record $m$, intervention family $\Pi$,
observable algebra $\mathcal A$, boundary and resolution declaration $B$,
resource declaration $C$, and horizon $\tau$, define
\begin{equation}
 \Efamily(m;\Pi,\mathcal A,B,C)
 =\left\{P^{\pi}_{\mathcal A}
   \bigl(F_{t:t+\tau}\mid m,B,C\bigr):\pi\in\Pi\right\}.
 \label{eq:future-family}
\end{equation}
The corresponding reachable support for policy $\pi$ is
\begin{equation}
 \Rset^{\pi}_{\tau}(m)=
 \operatorname{supp}P^{\pi}_{\mathcal A}
 \bigl(F_{t:t+\tau}\mid m,B,C\bigr),
 \label{eq:reachable-support}
\end{equation}
with topology and tolerance specified for continuous spaces. A single
continuous trajectory may have probability zero without being impossible.

Equation~\eqref{eq:future-family} is not itself an ICP law. It is shared
comparison language. A \emph{constraint} requires a matched contrast between
families, for example a natural history-formed carrier $m$ and a reference
$m_0$ prepared with the same declared ordinary state variables and resources:
\begin{itemize}
  \item a \textbf{hard constraint} changes reachable support on a common path
  or outcome space;
  \item a \textbf{soft constraint} changes weights while support may remain
  fixed; and
  \item an \textbf{operational constraint} changes a reachability or
  convertibility preorder under the same allowed operations.
\end{itemize}
No universal sign for an entropy difference follows from these definitions.

\subsection{A relative-entropy constraint diagnostic}

For a declared comparison $m$ versus $m_0$ at $x$, one distribution-sensitive member
of the future-family comparison is
\begin{equation}
 \mathfrak C_{x,\tau}^{\pi}
 (m\Vert m_0;\mathcal A,B,C)
 :=D_{\mathrm{KL}}\!\left(
 P^{\pi}_{\mathcal A}(F_{t:t+\tau}\mid m,B,C)
 \middle\Vert
 P^{\pi}_{\mathcal A}(F_{t:t+\tau}\mid m_0,B,C)
 \right),
 \label{eq:icp-constraint-divergence}
\end{equation}
provided both laws use the same outcome space, intervention, horizon, measure,
and coarse-graining, with the first law absolutely continuous with respect to
the second \citep{KullbackLeibler1951Information}. This is an indexed diagnostic,
not the definition of ICP: the full family in Eq.~\eqref{eq:future-family}
retains outcome identity, topology, interventions, and operational order that a
single scalar discards.

The IOTA support quantity has an exact specialization. If the post-intervention
law is uniform on
$\Rset_x^{\mathrm{post}}\subseteq\Rset_x^{\mathrm{prior}}$ and the reference
law is uniform on $\Rset_x^{\mathrm{prior}}$, then
\begin{align}
 D_{\mathrm{KL}}(U_x^{\mathrm{post}}\Vert U_x^{\mathrm{prior}})
 &=\sum_{f\in\Rset_x^{\mathrm{post}}}
   \frac{1}{|\Rset_x^{\mathrm{post}}|}
   \log\frac{|\Rset_x^{\mathrm{prior}}|}
                 {|\Rset_x^{\mathrm{post}}|} \notag\\
 &=\log\frac{|\Rset_x^{\mathrm{prior}}|}
                  {|\Rset_x^{\mathrm{post}}|}
 =\Delta I_x .
 \label{eq:iota-kl-identity}
\end{align}
This direction is essential: reversing the divergence is infinite under strict
support pruning. Non-nested or support-creating interventions require an
explicitly chosen alternative treatment rather than silent reuse of
Eq.~\eqref{eq:iota-kl-identity} \citep{Clark2026IotaConstraint}.

\subsection{The adjacent frontier and signed possibility-space flux}

The adjacent possible becomes operational only after its outcome types and
notion of ``one step'' are declared. Let
$r_\ell:\Omega\rightarrow\mathcal K_\ell$ be a frozen scale-$\ell$ map from
micro-outcomes to registered types or roles, and let $\mathcal V_t^\ell$ be the
types actualized by time $t$. Define
\begin{equation}
 \mathcal J_{t,\tau}^{\pi,\ell}(m_t)
 :=r_\ell\!\left(\Rset_\tau^\pi(m_t)\right)
   \setminus\mathcal V_t^\ell .
 \label{eq:adjacent-frontier}
\end{equation}
When $\tau$ is one registered transition, this is the novelty-bearing sector
of Kauffman's adjacent possible: types not yet actual but reachable in one step
from the actual \citep[p.~142]{Kauffman2000Investigations}. For longer horizons
it is an indexed accessible-novelty set, not a claim of literal one-step
adjacency. In either case it depends on $r_\ell,\pi,\tau,B,$ and $C$; changing a
classifier, resource budget, or policy can change the set.

For matched pre/post supports $R_0$ and $R_1$, retain the identity of the
change rather than only its cardinality:
\begin{equation}
 K=R_0\cap R_1,\qquad
 E=R_1\setminus R_0,\qquad
 L=R_0\setminus R_1,
 \label{eq:possibility-flux-ledger}
\end{equation}
where $E,L,$ and $K$ are enabled, lost, and retained possibilities. A signed
log-volume deformation is
\begin{equation}
 \chi(R_1:R_0):=\log\frac{\mu(R_0)}{\mu(R_1)}.
 \label{eq:signed-support-deformation}
\end{equation}
For nonempty finite uniform nested supports this gives the exact dual pair
\begin{equation}
 \chi(R_1:R_0)=
 \begin{cases}
 D_{\mathrm{KL}}(U_{R_1}\Vert U_{R_0}),&R_1\subseteq R_0
 \quad\text{(pruning)},\\[3pt]
 -D_{\mathrm{KL}}(U_{R_0}\Vert U_{R_1}),&R_0\subseteq R_1
 \quad\text{(enabling)}.
 \end{cases}
 \label{eq:signed-kl-duality}
\end{equation}
Thus $\chi>0$ denotes net contraction and $\chi<0$ net expansion under the
declared measure. If supports are non-nested, $\chi$ is only a signed net-volume
contrast: it can vanish under complete identity turnover. Report
Eq.~\eqref{eq:possibility-flux-ledger} and, where a scalar turnover diagnostic
is needed, a bounded common-support quantity such as
$D_{\mathrm{JS}}(U_{R_1},U_{R_0})$. Empty sets require explicit set reporting
or a declared absorbing outcome; they do not justify arbitrary pseudocounts.

\subsection{The persistent-carrier criterion}

A proposed informational carrier $M$ passes the current ICP criterion only if
four gates are met:
\begin{description}
  \item[Formation.] Controlled histories create the candidate distinction.
  \item[Sufficiency.] The distinction is sufficient for the declared future-law
  family to a frozen error tolerance.
  \item[Carrier intervention.] Ablating, replacing, transplanting, or scrambling
  the proposed carrier changes a preregistered component of the future family,
  while matched shams and resource controls do not.
  \item[Robustness.] The result survives adequate state augmentation and
  reasonable representation, partition, and coarse-graining changes, or its
  representation dependence is stated as a boundary.
\end{description}

Persistence across carrier succession adds one further demand. Let
$c_0,c_1,\ldots$ be physical realizations at successive times. They instantiate
the same role only relative to a declared equivalence relation $\sim_{\Pi,
\mathcal A,B,C,\tau}$ when replacement preserves the relevant future-response
family or a registered relational invariant. Material disjointness alone does
not establish sameness of role; sameness of an analyst's label does not
establish physical transport.

\subsection{Novelty and failure tests}

The strongest null is the complete-state ordinary-dynamics account: once every
causally relevant variable, law, input, boundary, and control is included,
there is no additional historical variable. State augmentation, predictive
causal states, process tensors, control/viability models, and
counterfactual-semantic measures are field-specific versions of this challenge.

ICP earns distinct physical content only by delivering at least one of the
following beyond the best native model:
\begin{enumerate}
  \item a minimal sufficient physically embodied equivalence class;
  \item a representation-stable carrier-succession invariant;
  \item a new impossibility or conversion theorem;
  \item a demonstrably superior compression of intervention-response families;
  or
  \item a novel, preregistered physical prediction.
\end{enumerate}
Repeated disappearance under adequate state augmentation, failure of matched
carrier interventions, dependence on one convenient representation, or the
need to import the target quantum structure would lower or defeat the relevant
claim.

\section{What the project has demonstrated}

Project evidence matters because IOTA is not only an interpretation of the
literature; it is an attempt to construct and intervene on the proposed
mechanism. Its evidential role must nevertheless remain bounded. A simulation
can establish consequences of its frozen rules and the realizability of a
formal mechanism. It cannot, without an independently validated bridge,
establish that nature uses those rules.

\subsection{Constraint transport and succession in IOTA}

The foundational IOTA result concerns a relational role rather than transport
of a marked payload \citep{Clark2026IotaConstraint}. Thirty-six previously
identified bridge-relocation events were replayed; twelve had a role-bearing
predecessor bridge, and three events across two seeds passed the strict test. In
each positive, tracked regions first lost connectivity and then regained it
through a causally necessary successor bridge with no cells in common with the
predecessor bridge. Removing that successor destroyed the restored relation.
Payload remained zero and the declared conservation ledger was exact.

The result supports a narrow statement: within this event class, a connectivity
constraint can be re-instantiated through complete replacement of its local
bridge. It does not show indefinite persistence, general carrier independence,
or correspondence with a particle or field. Follow-up tracking found finite
address-based relation lifetimes of 9, 3, and 51 recorded offsets. A prospective
expression-lineage rule extended two relations past original-address death to
170 offsets and at least 601 offsets; the third did not propagate.

Event-resolved interventions then separated route availability from
realization. Permitting a selected failed gate rescued six of six same-source
gate cases and eight of eight cases in a prospectively screened independent
geometry. Preventing source exhaustion reopened all 30 selected same-source
routes, but only 28 eventually handed off; an independent exhaustion case also
reopened the route without final rescue. An available future was therefore
causally relevant but not sufficient. Timing and competing events could close
the opportunity before realization.

Necessity tests revealed hidden redundancy. Among 164 factual branches with one
observed child, vetoing that child produced an alternate in 122 branches and no
alternate in 42. Sequentially suppressing newly realized successors yielded an
observed policy cut in 143 branches by depth eight; 19 remained right-censored
above that depth and two retained the original source at the horizon. An
independent cohort of 79 branches across four new source geometries produced 68
observed cuts, ten right-censored continuations, and one source-survival
censoring. These are adaptive, history-changing policy cuts, not a static count
of pre-existing interchangeable children.

These veto results have a bounded adjacent-possible interpretation. The first
veto exposed a different next realized type in 122 of 164 branches, and ordered
vetoes repeatedly changed which successor became available next. This is direct
project evidence for history- and policy-dependent exposure of local
alternatives, close to an adaptive version of
Eq.~\eqref{eq:adjacent-frontier}. It is not an enumeration of the complete
frontier. Branches were selected around an observed successor, some outcomes
were censored at the registered horizon, and all candidate states remained
inside a finite simulator with frozen update law and analyst-declared types.
The result therefore establishes neither open-ended evolution nor creation of
unprestatable physical state-space dimensions.

The evidence record also contains informative limits. Exact precursor placement
was causally necessary for later articulation in 25 selected channels, with
state-memory dependence in seven, but vetoing the first rebridge never sustained
the preregistered 128-offset separation. A proposed rooted graph-role signature
showed a directional contrast yet failed its full confirmation criteria
($p=0.0551$ for the primary paired test and a decisive failure of the separate
natural-closeness gate). A subsequent marked topology-by-state instrument was
feasible in development, but the independent cohort retained only 16 channels
against 20 required; no outcome was opened. The proper statuses are negative
for the frozen graph-role representation and untested for the independent
repair-affordance hypothesis.

Together, these results demonstrate beginnings of the proposed mechanism:
relational constraint replacement, lineage succession, finite opportunity,
causal realization, and temporal redundancy occur under frozen local rules and
prospective interventions. They do not demonstrate a substrate-independent
law.

The revised IOTA formalism also makes one exact mathematical connection. For
uniform laws on nested viable sets, its support log-ratio $\Delta I_x$ is
precisely the forward relative entropy in
Eq.~\eqref{eq:iota-kl-identity}. This upgrades a verbal entropy resemblance to a
proved identity inside the declared finite model. It does not show that IOTA's
viable-set measure is geometric, that nature realizes the model, or that the
scalar captures cardinality-preserving changes in future-law weights.

Equation~\eqref{eq:signed-kl-duality} extends that algebraic identity to nested
support expansion by reversing the divergence and retaining the sign. The
enabled/lost/retained ledger remains mandatory when pruning and enabling occur
together. This is a formal bridge supplied by the synthesis, not a new result
of the frozen IOTA experiments.

\subsection{Bounded Dirac role recovery}

The quantum-Dirac project asked whether a declared one-dimensional Dirac
operational role could be recovered across two separately registered local
quantum updates \citep{Clark2026DiracRole}. An ordered product walk and a
symmetry-derived direct quantum cellular automaton were tested with frozen
input sectors, mass and potential interventions, refinement ladders, observable
witnesses, causal cones, symmetry checks, resource counts, single-use
confirmation, and algorithmically separate numerical reconstruction.

The intact evidence supports qualified constructive sufficiency for the shared
bounded role core. A hostile post-result audit found that both primary and
replay implementations used the same factor-of-two error in an empirical
zone-edge predicate, selecting an empty sector. Those values and pass flags were
withdrawn; the independently derived analytic tail bound used the correct
sector. The audit also established that the free one-excitation blocks are
exactly related by a half-cell time-frame conjugacy after mass-angle remapping,
and that the reported resource difference depends on nonminimal gate-fusion
conventions. Thus the result does not establish two algebraically independent
architectures or a hardware advantage.

Most importantly for the emergence ladder, the construction supplies a quantum
Hilbert space, Born readout, clock, target map, intervention class, and mass
schedule. It shows that explicit local quantum mechanisms can realize the
bounded role; it does not yet show an IOTA coherence becoming a quantum
quasiparticle or mass emerging as an invariant. This is a successful rung and a
precise boundary, not a completed derivation.

\subsection{PSA as project evidence}

The submitted PSA semantic-field study constructs text signatures from
relations to a fixed catalog of propositional anchors
\citep{Clark2026PSA}. In one embedding model, 115-dimensional signatures retain
substantial neighborhood structure from a 4096-dimensional space. Disjoint
designed anchor catalogs produce correlated distance geometry beyond an
error-independence floor; random projections preserve more raw distance
structure but not the same shared distortion. A calibrated projection supports
auditable semantic blocking, while negative controls show that proximity does
not establish meaning equivalence.

PSA is relevant to ICP because it demonstrates a relational, frame-indexed
measurement architecture, tests cross-frame stability, and exposes how a
representation can constrain accessible comparisons. It is not empirical
fundamental physics. Its objects are sentences, embeddings, anchors, and an
engineered retrieval task. The evidence supports a methods analogy and a
project example of relational sufficiency and representation control; it does
not support the claim that the physical universe is semantic or that the
embedding space is an IOTA substrate.

This boundary also corrects an earlier Zipf interpretation. A hot-cell ceiling
in an adaptive voxel pipeline was removed by common-level reconstruction; a
pipeline-induced count feature cannot be promoted into a mechanism. The PSA
lesson retained for ICP is methodological: define the representation and
selection process before interpreting a regularity.

\section{Closest native formalisms and the redundancy test}

The most important literature for ICP is not the literature that sounds most
similar. It is the literature capable of making the proposal redundant. Five
clusters cover much of the generic future-constraint language: control and
viability, the adjacent possible, empowerment and causal information,
constructor and predictive-state theory, and feedback information
thermodynamics.

\subsection{Reachability and viability already describe constrained futures}

Hamilton--Jacobi reachability represents controlled future sets directly.
Mitchell, Bayen, and Tomlin characterize a finite-horizon backward-reachable set
for a nonlinear differential game as a sublevel set of a viscosity solution to
a Hamilton--Jacobi--Isaacs equation \citep{MitchellEtAl2005Reachable}. Viability
theory identifies states from which at least one admissible evolution can remain
inside a declared constraint set, with convergent approximation schemes under
stated conditions \citep{SaintPierre1994Viability}. These are not loose
analogies to Eq.~\eqref{eq:future-family}; they are mature mathematical theories
of accessible and survivable futures under controls and disturbances.

The distinction available to ICP lies upstream of the computed set. A
reachability calculation is given the state, law, targets, controls,
disturbances, boundary, and horizon. It need not ask how a particular
constraint role was formed, where it is physically realized, or whether it can
survive replacement of its first witness. If ICP stops at ``the state changes
what is reachable,'' it is standard control language. If it can identify a
minimal role that is formed and transported across realizations, it may add a
new state-equivalence or mechanism question to that language.

\subsection{The adjacent possible is the closest support-creation prior art}

Kauffman defines the chemical adjacent possible as molecular species absent
from the actual but synthesizable from it in one reaction step
\citep[p.~142]{Kauffman2000Investigations}. In the notation of
Eq.~\eqref{eq:adjacent-frontier}, this is a novelty-filtered, type-projected
one-step reachable set. That overlap is unusually strong: the adjacent possible
is native prior art for the part of ICP concerned not only with excluding
futures but with making new outcome types accessible.
Kauffman's later BioSystems formulation makes the stronger claim that enabling
constraints create unprestatable adjacent possibilities in biological
evolution \citep{Kauffman2014Prolegomenon}; a recent extension contrasts fixed
law-governed possibility spaces with an indefinite ``domain of no law''
\citep{PatraKauffman2025NoLaw}. These are important ontology-level competitors,
not empirically established consequences of the one-step definition.

The closest geometric competitor is still more direct. Fari\~na and Vallejo
define a ``biological cosmological constant'' $\Lambda_B$ proportional to the
squared fractional expansion rate of Fisher-information volume, interpret the
origin of life as a geometric phase transition, and apply the construction to
staged Phanerozoic diversification \citep{FarinaVallejo2026LambdaB}. This is
explicit prior art for geometric expansion into the adjacent possible, not
evidence for ICP. Its eco-evolutionary manifold and nonnegative squared rate do
not identify a replaceable carrier, preserve enabled/lost/retained identities,
or distinguish signed local pruning from enabling. ICP therefore cannot claim
the first geometric adjacent-possible measure or a generic phase-transition
signature; its narrower opportunity is an intervention-defined signed contrast
and a prospectively predicted carrier-linked crossover.

Later work turns the idea into several non-equivalent models. The urn model
with triggering adds $\nu+1$ previously unseen colors when a novel color is
drawn, while reinforcement $\rho$ increases reuse; it yields correlated
novelty together with Heaps- and Zipf-like scaling in the stated regimes and
was compared with textual, music, tagging, and Wikipedia sequences
\citep{TriaEtAl2014Novelties}. Edge-reinforced walks instead place discovery on
a declared concept network \citep{IacopiniEtAl2018Innovation}; multi-agent urns
model evolving social contacts \citep{UbaldiEtAl2021SocialNetworks}; and
higher-order novelty counts first occurrences of relations or combinations,
which can distinguish sequences with similar first-order discovery rates
\citep{DiBonaEtAl2025HigherOrder}. Dynamic scaffolding can also change which
possibilities are developmentally available, although the published treatment
is a modeling proposal rather than a direct test
\citep{Griesemer2024Scaffolded}. Numerical work in a Kauffman chemical model
further reports that autocatalytic closure can causally change complexity
dynamics under declared interventions \citep{Faggian2025Closure}; causal
organization-to-novelty effects are therefore prior art as well as conceptual
expectations.

The combinatorial TAP family makes the expansion explicit. One studied form is
\begin{equation}
 M_{t+1}=(1-\mu)M_t+
 \sum_{i=2}^{M_t}\alpha_i\binom{M_t}{i},
 \label{eq:tap-equation}
\end{equation}
where subsets of $M_t$ existing types generate candidates and $\mu$ represents
loss. Stochastic and deterministic variants can produce a long plateau followed
by rapid or finite-time explosive growth
\citep{SteelEtAl2020Combinatorial,CortesEtAl2025TAP}. This behavior is not a
universal observational law. Explicit-type models show that duplicate
rediscovery and finite worlds matter \citep{VillaniSerra2023BinaryString}, and a
1908--2016 Swedish product study found resource-constrained innovation rates,
Heaps-like type discovery, and useful but incomplete predictive information in
product-space topology rather than unrestricted combinatorial blow-up
\citep{Taalbi2026AdjacentPossible}.

This literature changes the novelty verdict. ICP cannot claim invention of
history-dependent accessibility, triggered support growth, novelty clustering,
scaffolded opening of opportunities, or a generic plateau-to-explosion
transition. It also cannot claim the first Fisher-geometric or cosmological
formalization of adjacent-possible expansion. Nor do Heaps or Zipf exponents
identify the mechanism. The possible
remaining addition is causal and physical: predeclare the type map, frontier,
carrier, horizon, and resources; alter or replace that carrier; and predict the
enabled, lost, and retained frontier on held-out interventions. Formal AP models
usually prestate an alphabet, graph, urn, or combinatorial generator. The
stronger claim that biological phase space is intrinsically unprestatable
\citep{LongoEtAl2012Enablement,Kauffman2014Prolegomenon,
PatraKauffman2025NoLaw} cannot simultaneously serve as a directly enumerated
experimental variable without an operational surrogate.

\subsection{Empowerment and future-option measures}

Empowerment makes the overlap still sharper. Its original finite-step form is
the channel capacity from present action sequences to later sensed states---a
scalar measure of potential perceivable control \citep{KlyubinEtAl2005Empowerment}.
Process empowerment replaces fixed open-loop scripts with feedback-conditioned
probing policies and can expose controllability hidden by noise or traps
\citep{TiomkinEtAl2025ProcessEmpowerment}. Causal entropic forcing instead
builds a finite-horizon path-entropy objective into an effective force and
produces adaptive-looking behavior in simple simulations
\citep{WissnerGrossFreer2013CausalEntropy}.

These works defeat any claim that ICP discovered the idea of valuing or
counting future options. They also clarify what ICP does \emph{not} assert. The
principle does not presently say that systems maximize reachable-set size,
empowerment, entropy, or optionality. A constraint may reduce, redirect, or
reorganize futures; persistent structure can increase viable options at one
scale while excluding paths at another. Scalar option values also collapse
outcome identity: equal capacities need not imply equal future-law families.

\subsection{Causal information and the closest semantic rival}

Causal-information measures require post-intervention distributions or a
defended causal model. Ay and Polani formalize information flow using causal
rather than merely observational independence \citep{AyPolani2008InformationFlows}.
Counterexamples show why transfer entropy and related dyadic quantities can
both overestimate flow and underestimate influence
\citep{JamesEtAl2016InformationFlows}. Predictive association, including a
Zipf law or high mutual information, is therefore not a carrier test.

Kolchinsky and Wolpert provide the closest mature formal rival to the phrase
``information that makes a difference.'' They define semantic information as
system--environment correlation whose controlled scrambling lowers a declared
future viability \citep{KolchinskyWolpert2018Semantic}. The intervention,
system boundary, ensemble, horizon, syntactic measure, and viability function
are constitutive. On a survival or self-maintenance task, an ICP analysis should
specialize or compare with this framework, not silently rename it.

ICP is broader only in a limited sense: its target may be any registered
future-response role, not necessarily viability, and it emphasizes formation
and succession across original-witness replacement. Breadth is not novelty by
itself. Unless carrier succession produces a new invariant or empirical
prediction, semantic information and ordinary causal-effect analysis may
already exhaust the phenomenon.

\subsection{Constructor theory and predictive-state theory}

Constructor theory supplies a law-level language of possible and impossible
repeatable tasks and characterizes information variables through physical task
repertoires such as copying, permutation, and distinguishability
\citep{Deutsch2013Constructor,DeutschMarletto2015Information}. This is the
cleanest neighbor for hard physical possibility. It does not, in its basic task
definition, supply a history-conditioned stochastic process or require that a
present constructor was formed through carrier succession. ICP may eventually
add a temporal formation and inheritance layer; for now the relationship is
structural, not derivational.

Computational mechanics begins from histories rather than tasks. Causal states
are equivalence classes of pasts with the same conditional distribution over
futures and yield a minimal optimal predictive representation under the stated
process assumptions \citep{ShaliziCrutchfield2001ComputationalMechanics}.
Input--output epsilon-transducers extend the construction to conditional output
laws under future input sequences \citep{BarnettCrutchfield2015Transducer}.
This is almost exactly the $H\rightarrow M\rightarrow P(F\mid M)$ part of
ICP.

The difference is physical, not algebraic. An inferred causal state need not be
localized, manipulable, or Pearl-causal; conditioning on an input sequence is
not automatically a surgical intervention. ICP's carrier criterion asks
whether a physically embodied realization is necessary and sufficient for the
relevant family and whether its replacement preserves the role. If that
physical test adds nothing to the predictive equivalence class, computational
mechanics is the more mature account.

\subsection{Feedback information thermodynamics}

Feedback thermodynamics gives direct evidence that a retained measurement can
change a later physical transformation. Sagawa and Ueda's generalized
Jarzynski relation includes information and feedback efficacy in work bounds
\citep{SagawaUeda2010Feedback}; Toyabe and colleagues tested information-to-free-
energy conversion with position-conditioned feedback on a Brownian particle
\citep{ToyabeEtAl2010InformationEnergy}. An autonomous bit-tape demon makes the
memory carrier and its thermodynamic accounting explicit
\citep{MandalJarzynski2012Demon}. Current work also shows that finite memory and
a limited protocol repertoire alter the optimal feedback value
\citep{HartleEtAl2026LimitedResources}.

A 2026 experimental preprint provides an especially relevant boundary case.
Time-delayed measurements of a colloid in a viscoelastic bath reveal memory not
captured by the visible instantaneous coordinate, and feedback can extract work
from hidden environmental degrees of freedom \citep{MurugaEtAl2026HiddenDegrees}.
This is strong evidence that information can persist outside the original or
currently observed bearer while remaining physically realized elsewhere. It
also supports the state-augmentation null: the enlarged system is Markovian and
the hidden modes carry ordinary excess free energy. As a preprint, it remains
frontier evidence.

Information thermodynamics therefore supports causal utility and closed
physical accounting. It does not support the inference that energy emerges from
information, because energy, Hamiltonians, baths, memories, and protocols enter
the theory as supplied structure. IOTA must derive mass- and energy-like
invariants at a later rung rather than cite feedback conversion as that
derivation.

\subsection{Result of the redundancy audit}

The generic ICP proposition is largely preempted. Mature theories already
establish that present physical states, controls, memories, and correlations
differentiate future possibilities and performance. The viable novelty target
is narrower:

\begin{quote}
An interaction forms a physically embodied constraint role; the original
witness ceases to instantiate it; a successor realization preserves the
registered intervention-response role; and perturbing that successor changes
the future family under matched controls.
\end{quote}

That target aligns with the strongest IOTA evidence and remains open as a
general principle. It is also falsifiable by complete-state reduction,
failure of role equivalence, or representation dependence.

\section{Quantum records and the emergence bridge}

Quantum theory supplies the sharpest positive precedents for information
leaving its first bearer, and the sharpest warning against treating that
precedent as a derivation of quantum theory.  Once a Hilbert space, a tensor or
algebraic decomposition, a state, and an intervention class are supplied,
standard quantum mechanics already explains how interactions distribute
correlations, how some observables become stable under environmental
monitoring, and how collective excitations acquire particle-like roles.  The
IOTA hypothesis begins one level earlier.  It must explain how its local
informational distinctions could acquire the composition, phase, probability,
and transformation structure presupposed by those results.

Table~\ref{tab:emergence-ladder} makes this separation explicit.  The first
rung, and precursor aspects of the second, have bounded IOTA evidence; no
project result yet establishes quantum coherence.  The current Dirac
construction occupies a higher, supplied-quantum rung.  The remaining arrows
are research hypotheses, not consequences of similarity in vocabulary.

\begin{table}[t]
\centering
\small
\caption{The emergence ladder and the evidential status of each bridge.}
\label{tab:emergence-ladder}
\begin{tabularx}{\linewidth}{@{}>{\raggedright\arraybackslash}p{0.16\linewidth}
  >{\raggedright\arraybackslash}X >{\raggedright\arraybackslash}X@{}}
\toprule
Rung & Present support & Missing discriminating result \\
\midrule
Local interaction $\rightarrow$ transported constraint
& Project construction: selected connectivity roles are re-instantiated in
disjoint successors, with causal intervention gates.
& Replication across rules, schedules, topologies, and physically equivalent
representations. \\
Transport $\rightarrow$ persistent coherence
& Project lineages sometimes outlive original addresses, a pre-quantum
persistence result; quantum and biological records are independent precedents.
& A minimal compositional invariant whose persistence is not an artifact of a
chosen tracking rule. \\
Coherence $\rightarrow$ quantum structure
& No base-IOTA derivation.  Process tensors, decoherence, reference frames, and
the Dirac project all assume quantum theory.
& Derive state composition, reversible/interventional dynamics, interference,
and probability calibration without importing amplitudes or the Born rule. \\
Quantum structure $\rightarrow$ quasiparticle
& Established local quantum models and experiments realize fractionalized and
topological excitations.
& Create and manipulate an IOTA-derived excitation with declared dispersion,
statistics, symmetry, conservation, and scattering. \\
Quasiparticle $\rightarrow$ fields, mass, and energy
& One-dimensional quantum cellular automata recover bounded free-field roles;
the project recovers a bounded one-particle Dirac role.
& Interacting multiparticle limits, no-go closure, and derived rather than
scheduled mass- and energy-like invariants. \\
\bottomrule
\end{tabularx}
\end{table}

\subsection{Multi-time quantum memory is a native future-law formalism}

\paragraph{Established field result.}
The process-tensor framework represents the statistics of a quantum process
under a sequence of possible interventions.  It separates the experimenter's
operations from the multi-time object that maps those operations to outcome
probabilities; quantum Markov order and memory classifications then depend on
the chosen instruments and system--environment split
\citep{PollockEtAl2018ProcessTensor,PollockEtAl2018Markov,
TarantoEtAl2019MarkovOrder,TarantoEtAl2024ClassicalMemory}.  Recent process
tomography demonstrates that such objects can be reconstructed in controlled
devices, subject to finite-dimensional, sampling, and control assumptions
\citep{GiarmatziEtAl2025Tomography}.  This is close to
Eq.~\eqref{eq:future-family}: the future object is not one unconditioned state
but a family of responses to allowed operations.

\paragraph{Synthesis inference.}
A multi-time process is a natural quantum benchmark for an ICP carrier claim.
Two histories can leave the same measured reduced system state yet differ in
environmental correlations and therefore in their later intervention-response
families.  The record has not escaped physics; it resides in an enlarged joint
state or process.  Consequently, state enlargement is both a positive method
for locating memory and the strongest null against a new historical variable.

\paragraph{ICP hypothesis.}
The stronger claim is that a compact, physically identifiable role can be
formed, edited, and transferred while preserving the relevant process family.
That claim would require histories matched on the best available present state,
randomized preservation or rewriting of the proposed carrier, resource-matched
shams, and held-out multi-time interventions.  If the augmented standard
process tensor explains every effect without a new invariant or compression,
ICP remains a useful organization of ordinary quantum memory rather than an
additional quantum law.

\subsection{Environmental records: a real mechanism with an objectivity ladder}

\paragraph{Established field result.}
Decoherence disperses phase information into correlations with an environment,
and environment-induced superselection identifies observables comparatively
stable under the specified interaction.  Quantum Darwinism adds a distinct
condition: many disjoint fragments may carry information about a common
preferred observable.  Environment-as-witness constructions and a broad
finite-dimensional channel theorem show how fragment channels can become
approximately measure-and-prepare channels for one common POVM
\citep{OllivierPoulinZurek2004Witness,
BrandaoEtAl2015Darwinism}.  Controlled photonic and spin experiments
have observed redundant imprints under engineered or well-characterized
conditions \citep{CiampiniEtAl2018DarwinismExperiment,ChenEtAl2019DarwinismExperiment,
UndenEtAl2019Classicality}.

This is direct evidence that an interaction-generated distinction can be
embodied in secondary degrees of freedom and read away from the source.  It is
not a demonstration of information without a carrier.  Nor do the familiar
levels collapse into one another:
\begin{equation}
 \text{decoherence}\not\equiv\text{pointer selection}
 \not\equiv\text{redundant encoding}
 \not\equiv\text{accessible consensus}
 \not\equiv\text{objective broadcast record}.
 \label{eq:objectivity-ladder}
\end{equation}
Spectrum broadcast structure supplies a stringent state-level form in which
multiple fragments have distinguishable conditional supports.  Strong quantum
Darwinism additionally requires locally accessible classical information, and
the multipartite-independence condition must be stronger than pairwise
conditional independence \citep{HorodeckiEtAl2015Objectivity,
LeOlayaCastro2018Objectivity,FellerEtAl2021CommentDarwinism}.  Mutual-information
plateaus, redundancy, and observer consensus can therefore disagree
\citep{ChisholmEtAl2023RedundancyConsensus}.  Classification results further show that
record proliferation requires dynamical structure; arbitrary coupling does not
automatically make an objective record \citep{DoucetDeffner2024DarwinismHamiltonians}.

\paragraph{Synthesis inference.}
Quantum Darwinism supplies the clearest independent model of persistence beyond
an original bearer: correlations generated at a source can remain available in
other physical carriers.  It does not yet supply the ICP succession criterion.
Most experiments ask what can be inferred from a fragment now, not whether
erasing, swapping, or replacing that fragment later changes a registered
future process while all ordinary resources are controlled.

\paragraph{Factorization boundary.}
The words ``system,'' ``environment,'' and ``fragment'' presuppose a
decomposition.  The same global state can display redundant imprints of
noncommuting observables under different, non-comparable environment
partitions \citep{Fu2021RedundantImprints}; a fixed Hamiltonian can also admit
multiple tensor-product structures with robust pointer states that need not
track familiar locality \citep{AdilEtAl2026ClassicalSubsystems}.  This does not
undermine an operationally fixed laboratory partition.  It blocks the stronger
inference from a redundant record to a unique, observer-independent carrier or
classical world.  An ICP quantum claim must therefore state the factorization
or observable algebra, accessible fragment operations, reference resources,
and the transformations under which carrier identity is expected to survive.

\subsection{Relational facts and quantum reference frames are not memories}

Relational quantum mechanics (RQM) and quantum reference-frame (QRF) theory
share an insistence that descriptions are relational, but they answer different
questions.  RQM is an interpretation in which variable values are facts
relative to interactions; later work distinguishes approximately stable facts
from generic relative facts \citep{Rovelli1996RelationalQuantum,
DiBiagioRovelli2021StableFacts}.  Whether its commitments permit consistent
cross-perspective fact assignments remains actively disputed
\citep{Pienaar2021Quandaries,DiBiagioRovelli2022Reply}.  That dispute is
interpretive and supplies no independent experimental confirmation of IOTA.

QRF formalisms instead transform states, observables, dynamics, and
measurements relative to physical quantum systems.  Superposition and
entanglement can be frame-relative while consistently transformed operational
probabilities remain invariant \citep{GiacominiEtAl2019QuantumFrames}.  Operational
equivalence and reversibility depend on the framework, localization, symmetry
constraints, and access to global charges
\citep{CaretteEtAl2025OperationalFrames,DoatVanrietvelde2026Charge}.

\paragraph{Synthesis inference.}
This literature requires ICP to distinguish a changed description from a
changed future.  A frame transformation may relabel subsystems and
entanglement without changing any invariant probability.  A relative fact may
be created in an interaction without becoming a durable, manipulable memory.
Neither implication
\begin{equation}
 \text{relational fact}\Rightarrow\text{stored carrier},\qquad
 \text{frame-relative state}\Rightarrow\text{different reachable family}
\end{equation}
holds without an additional physical construction.  A proposed carrier must
either survive the relevant frame equivalences or be explicitly indexed to a
declared frame and observer algebra.

\subsection{The Dirac construction is a rung, not the missing arrow}

The bounded Dirac result belongs between supplied quantum structure and an
effective one-particle role.  Its local update specifications, convergence
tests, causal cones, interventions, and observable witnesses show how a known
role can be recovered constructively \citep{Clark2026DiracRole}.  The hostile
audit and exact time-frame relation limit the architecture-independence claim,
but do not erase the intact shared role core.

Its limitation is decisive for the present paper: complex amplitudes, Hilbert
composition, unitary updates, a Born readout, a physical clock, and a mass
schedule enter as supplied structure.  Thus the construction answers
``can these local quantum updates realize this bounded Dirac role?''  It does
not answer ``can IOTA interactions create quantum state structure?''  A visual
analogy between an IOTA coherence and a wave packet would not close this gap.
The required bridge must recover interference and composition under novel
tests, not encode the target equation in a relabeled update rule.

\subsection{Quasiparticles establish possibility, not provenance}

\paragraph{Established field result.}
Local quantum organization can produce excitations whose operational identity
is not that of any microscopic constituent.  Exactly soluble string-net phases
support emergent gauge and fermionic roles, and the Kitaev honeycomb model
supports fractionalized Majorana and flux excitations
\citep{LevinWen2005StringNet,Kitaev2006Anyons}.  Material measurements provide
model-dependent evidence for such fractional excitations, while controlled
processors and quantum gases have created and manipulated non-Abelian or
anyonized collective behavior
\citep{JansaEtAl2018Fractional,IqbalEtAl2024NonAbelian,
DharEtAl2025Anyonization}.  These are not merely metaphorical particles: within
their domains they have measurable statistics, braiding, correlations, or
response signatures.

One-dimensional quantum cellular automata provide an especially close
existence precedent.  A multiparticle QCA can match its quantum walk in the
one-particle sector and approach a free-fermion QFT at long wavelength.  The
same analysis proves an obstruction for the proposed class in higher
dimensions \citep{MlodinowBrun2020QCAQFT}.  The positive construction and the
no-go result are equally relevant: local updates can recover a field role, but
not generically and not without dimensional and structural conditions.

\paragraph{Synthesis inference.}
These results defeat the blanket objection that microscopic carrier variables
and effective particle identity must coincide.  They make the later portion of
the IOTA ladder logically plausible.  They do not identify the IOTA lattice as
nature's substrate because every cited example starts with a quantum Hilbert
space, Hamiltonian or circuit, selected phase, and ordinary laboratory
resources.

\paragraph{ICP hypothesis and gate.}
The next project object must be more than a long-lived pattern.  A
prospectively declared excitation should have creation and annihilation
operations, finite propagation, dispersion, composition and exchange rules,
symmetry and conserved quantities, scattering or perturbative response, and
robustness under carrier replacement and representation change.  It must then
survive the imported-quantum, lattice-artifact, encoding, and single-target
nulls.  Passing those gates would establish a quasiparticle-like role in the
model.  Only a prior derivation of the quantum structure would make it evidence
for the full substrate claim.

\subsection{Causal order is a boundary on the word ``future''}

The ordinary form of Eq.~\eqref{eq:future-family} presumes a declared temporal
or operational order.  Process-matrix theory and the quantum switch show that
quantum operations can be organized without one fixed global order, and causal
witness experiments have certified causal nonseparability under trusted device
models \citep{OreshkovEtAl2012CausalOrder}.  Circuit order, process
nonseparability, causal-inequality violation, and indefinite order among
localized spacetime events are nevertheless distinct claims
\citep{PaunkovicVojinovic2020CausalOrders,OrmrodEtAl2023SectorialConstraints}.
A 2026 photonic experiment reported a Bell-like causal-order-inequality
violation at 5.7 standard deviations under stated free-intervention and
relativistic-separation assumptions, while retaining locality, fair-sampling,
free-choice, and setup loopholes \citep{GuoEtAl2026CausalOrder}.

\paragraph{Boundary verdict.}
A quantum control can condition the family of later operation sequences.  That
control is not thereby a history-generated record; causal nonseparability is
not thereby indefinite spacetime geometry; and neither result derives
causality from information.  Where order is not fixed, ICP must replace an
ordinary time-indexed future with an explicitly ordered process functional.
The causal-order literature is therefore a formal boundary condition on the
proposal, not evidence for the substrate ontology.

\subsection{Quantum-sector conclusion}

The quantum literature supports three ingredients: interactions can create
physically separated records; record accessibility is relative to an algebra,
factorization, reference frame, and operation class; and organized quantum
substrates can realize genuine particle-like excitations.  The synthesis
inference is that persisted information should be treated as an embodied
source--carrier--operation relation rather than a free-floating substance.
The open IOTA/ICP hypothesis is that such roles can be constructed from the base
informational dynamics and can acquire quantum composition without the quantum
target being supplied.  No result reviewed here closes that arrow.

\section{Cross-domain stress tests: statistics, entropy, life, and gravity}

The following domains are not four independent confirmations of ICP.  They are
four ways to try to break it.  Zipf statistics test whether a compelling
regularity identifies a mechanism.  Entropy tests whether one scalar retains
the structure of constrained possibility.  Biology tests whether history can
build and maintain a causally effective record under ordinary material
turnover. Gravity tests whether local geometric entropy can be mapped to an
operational future constraint without choosing the operands for agreement;
black-hole information further tests whether ``preserved,'' ``outside,'' and
``recoverable'' remain equivalent when observables, factorization, and
computational resources become nontrivial. In each case,
the native theory carries the scientific result; ICP is at most a proposed
mapping until it supplies an additional invariant, prohibition, compression,
or prediction.

\subsection{Zipf regularity underdetermines its cause}

\paragraph{Established field result.}
A roughly inverse rank--frequency curve can be generated by incompatible
mechanisms.  Cumulative-advantage and innovation dynamics produce skew
distributions \citep{Simon1955Skew}; random character strings provide a
memory-free counterexample \citep{Li1992RandomTexts,
FerrerCanchoElvevag2010RandomTexts}; communicative least-effort models use a
signal--referent optimization \citep{FerrerCanchoSole2003LeastEffort}; broad
latent-variable mixtures can yield Zipf and apparently critical behavior
without intrinsic critical tuning \citep{SchwabEtAl2014LatentVariables,
AitchisonEtAl2016UnderlyingVariables}; and sample-space-reducing processes
generate the law through explicitly history-dependent successor pruning
\citep{CorominasMurtraEtAl2015SSR,ThurnerEtAl2015SentenceSSR}.  The same coarse
shape can therefore occur with no memory, hidden common causes, optimization,
or genuine path dependence.

Rank--count scaling, the distribution of type counts, and vocabulary growth
are also different statistical objects.  Their exponent relations are only
conditional or asymptotic, with logarithmic corrections near the canonical
Zipf case \citep{RosilloRodesEtAl2026TypeToken}.  Holding the final count
multiset fixed while permuting event order can produce different discovery
paths \citep{ZimmerlinEtAl2026Discovery}.  A second scaling plot is thus not an
independent causal witness if it is algebraically or procedurally coupled to
the first.

\paragraph{Strongest rival explanations.}
Finite samples, aggregation, censoring, adaptive preprocessing, topical
mixtures, static grammar or network structure, and transient vocabulary growth
can all imitate a proposed constraint mechanism.  Visual linearity or a high
$R^2$ is not a distributional test.  Likelihood-based fitting, absolute
goodness of fit, alternative heavy-tail comparisons, dependence-aware
uncertainty, and replication are prerequisites even for the modest claim that
a power law is supported on a declared range
\citep{ClausetEtAl2009PowerLaws,CorralEtAl2020DistinctFlavors}.

\paragraph{Synthesis inference.}
Zipf is useful to ICP precisely because it demonstrates causal
underdetermination.  The discriminating object is not an exponent but the law
of future reuse, novelty, and succession among histories matched on present
marginal counts.  Within-document shuffles, block shuffles, frozen prefixes,
latent conditioning, and held-out prediction can reveal path dependence while
preserving a marginal.  Those are analytic interventions; they do not yet
manipulate a physical carrier.

\paragraph{ICP hypothesis.}
A sequence-bearing role becomes ICP evidence only if a candidate embodied
record is formed, is sufficient for the declared future law, survives rival
conditioning, and changes that law when its carrier is erased, exchanged, or
replaced under matched controls.  If counts, latent variables, and ordinary
dynamics absorb the effect, the Zipf result rejects rather than supports the
strong claim.  No Zipf marginal in the present project establishes ICP.

\subsection{Constraint appears in entropy only after the problem is declared}

The phrase ``constraint lowers entropy'' is true in special nested problems and
false as a general identification.  If $C_2\subseteq C_1$ are measured by the
same $\mu$, then the Boltzmann-style support quantity obeys
\begin{equation}
 k\log\mu(C_2)\leq k\log\mu(C_1).
 \label{eq:nested-support}
\end{equation}
Adding a feasible constraint to the same maximum-entropy optimization likewise
cannot raise the optimized objective.  A physical intervention, however, can
also change dynamics, Hamiltonian, reference measure, observable algebra,
coarse-graining, operation class, or horizon.  Under those changes an entropy
may rise, fall, or remain fixed.

\paragraph{Established field results.}
Shannon entropy quantifies average coding uncertainty for a declared
distribution \citep{Shannon1948Communication}.  Boltzmann multiplicity and
Jaynesian maximum entropy use closely related mathematics for different
scientific jobs: compatible physical microstates versus inference subject to
selected expectation constraints \citep{Boltzmann1877Relationship,
Jaynes1957InformationTheory}.  Lieb and Yngvason begin instead with an
adiabatic-accessibility preorder and, under equilibrium axioms, derive an
entropy representing that order; outside that regime lower and upper entropies
can leave genuine incomparability \citep{LiebYngvason1999SecondLaw,
LiebYngvason2013Nonequilibrium}.  Here accessibility is logically prior to its
scalar representation.

Temporal organization separates one-symbol entropy, entropy rate, predictive
information, and predictive-state memory.  A fair independent coin and an
alternating binary process with random initial phase both have one bit of
single-symbol entropy, while only the latter stores a bit that fixes every
later symbol.  Entropy-convergence and computational-mechanics analyses make
this distinction formal \citep{CrutchfieldFeldman2003Regularities,
ShaliziCrutchfield2001ComputationalMechanics}.  Likewise, topological entropy
tracks growth of distinguishable path support, whereas metric entropy depends
on the invariant probability measure \citep{Kolmogorov1958MetricInvariant,
Sinai1959Entropy,AdlerEtAl1965TopologicalEntropy}.  Hard possibility and soft
weighting can therefore vary independently.

Path-level and thermodynamic quantities add other projections.  Large-
deviation rate functions retain the exponential cost of rare future sectors,
but conditioning on a rare future is not a physical intervention that causes
it \citep{DonskerVaradhan1975MarkovExpectations,
ChetriteTouchette2015Conditioned}.  Fluctuation relations compare forward and
reverse path probabilities and work statistics; they do not measure stored
predictive organization \citep{Jarzynski1997Equality,
Crooks1999EntropyProduction,Seifert2005TrajectoryEntropy}.  In a specified
driven stochastic model, nonpredictive information is tied to dissipation, a
powerful but conditional connection rather than a universal identity
\citep{StillEtAl2012Prediction}.

Quantum and finite-resource settings make scalar insufficiency sharper.
Conditional min- and max-entropies answer operational one-shot tasks at a
chosen error tolerance \citep{KonigEtAl2009MinMax}.  Pure-state LOCC
convertibility is governed by majorization rather than entanglement entropy
alone \citep{Nielsen1999EntanglementTransformations}.  Thermal resource
theories can require a family of generalized free energies, with additional
coherence constraints when superpositions across energies matter
\citep{BrandaoEtAl2015SecondLaws,LostaglioEtAl2015Coherence}.  Restricting
control Hamiltonians changes extractable work while the state remains fixed
\citep{WilmingEtAl2016ControlRestrictions}.  In gauge theories even the entropy
of a region depends on the chosen observable algebra, center, or edge-mode
prescription \citep{Donnelly2012GaugeEntropy,CasiniEtAl2014GaugeFields}.

\begin{table}[t]
\centering
\small
\caption{Constraint structure and entropy are not one-to-one.}
\label{tab:constraint-entropy}
\begin{tabularx}{\linewidth}{@{}p{0.23\linewidth}X X@{}}
\toprule
Declared constraint object & Native probe & Information lost by one scalar \\
\midrule
Compatible states or paths & Boltzmann/support or topological entropy
& identity of excluded states, graph structure, and intervention response \\
Probability weights & Shannon, R\'enyi, or metric entropy
& outcome identity, causal provenance, and rare-event geometry \\
Predictive process & entropy rate, excess entropy, statistical complexity
& full causal-state topology and physical carrier location \\
Forward/reverse experiment & entropy production and fluctuation relations
& stored memory and support accessibility \\
Allowed transformations & majorization and resource-monotone families
& incomparability and operations hidden by the chosen monotone \\
Observable algebra & von Neumann or algebraic entropy
& distinctions inaccessible outside the selected algebra \\
\bottomrule
\end{tabularx}
\end{table}

\paragraph{Synthesis inference.}
Constraint is upstream only in a carefully indexed sense: the compatible set,
path law, operation preorder, observable algebra, reference state, task, and
resource budget are specified first; an entropy is then a functional or
monotone of that structure.  Equal entropy does not imply equal accessible
futures, and lower entropy does not imply more organization.  Feedback
thermodynamics establishes that a recorded distinction can improve work
extraction within ordinary energetic accounting
\citep{SagawaUeda2010Feedback,ToyabeEtAl2010InformationEnergy}; it does not
derive energy from information.

\paragraph{ICP hypothesis.}
An embodied record may change support, weights, predictive memory, rare-event
cost, dissipation, or a conversion preorder.  The experiment must say which,
then compare the complete native model after carrier intervention and state
augmentation.  Searching across entropy definitions until one moves in the
desired direction would be representation selection, not evidence for a
principle.

\subsection{Gravity from Entropy supplies a geometric bridge target}

\paragraph{Established field result within the GfE construction.}
Gravity from Entropy (GfE) defines higher-order metric operators, assumes their
flattened ratio is positive, and writes its gravitational log-determinant term as
\begin{equation}
 \mathcal L_{\mathrm{GfE}}
 =-\operatorname{Tr}_{F}\ln\boldsymbol\Theta,
 \qquad
 \boldsymbol\Theta=\widetilde{\mathbf G}\,\widetilde g^{-1},
 \qquad
 \boldsymbol{\mathcal G}=\boldsymbol\Theta^{-1}
 \label{eq:gfe-logdet}
\end{equation}
\citep{Bianconi2025Gravity,Bianconi2026Thermodynamics}. These metric operators
are not normalized probability states, and the subscript $F$ denotes the
flattened-tensor trace rather than a probabilistic state space. Thus
Eq.~\eqref{eq:gfe-logdet} is not automatically a Kullback--Leibler or Umegaki
divergence. In the thermodynamic development, an eigenmode contributes
\begin{equation}
 c_k:=\frac{\delta s_k}{\delta v}=\ln\mathcal G_k,
 \qquad
 \theta_k=e^{c_k}-1,
 \qquad
 \frac{\delta\epsilon_k}{\delta v}=e^{c_k}-1-c_k.
 \label{eq:gfe-thermodynamics}
\end{equation}
Individual $c_k$ can be signed; a nonnegative mode-by-mode identification with
Eq.~\eqref{eq:icp-constraint-divergence} is therefore generally invalid. The
signed support deformation in Eq.~\eqref{eq:signed-support-deformation} removes
that algebraic sign obstruction, but does not by itself supply a physical mode
map.

\paragraph{Synthesis inference: the bridge formula.}
Let $Q_x^{\mathrm{prior}}$ and $Q_x^{\mathrm{post}}$ be as-yet-undefined positive
operators whose flattened determinants encode one declared measure of IOTA
reachable futures. The proposed bridge is
\begin{equation}
 \begin{aligned}
 \boldsymbol\Theta_x
 &\stackrel{?}{=}(Q_x^{\mathrm{prior}})^{-1/2}
 Q_x^{\mathrm{post}}(Q_x^{\mathrm{prior}})^{-1/2}
 \sim Q_x^{\mathrm{post}}(Q_x^{\mathrm{prior}})^{-1},\\
 \det_F Q_x&=a\,\mu(\Rset_x),
 \qquad a>0\ \text{common to prior and post}.
 \end{aligned}
 \label{eq:gfe-iota-operator-bridge}
\end{equation}
Here $\sim$ denotes similarity; the first representative is positive and has
the determinant ratio of the original product. For the finite IOTA identity,
$\mu$ is counting measure. Under these algebraic conditions,
\begin{equation}
 \begin{aligned}
 -\ln\det_F\boldsymbol\Theta_x
 &=\ln\frac{\mu(\Rset_x^{\mathrm{prior}})}
             {\mu(\Rset_x^{\mathrm{post}})}
 =:\chi_x.
 \end{aligned}
 \label{eq:gfe-iota-scalar-bridge}
\end{equation}
For finite uniform nested supports, Eq.~\eqref{eq:signed-kl-duality} gives the
orientation explicitly: $\chi_x$ is the forward KL divergence and equals
IOTA's $\Delta I_x$ under pruning, whereas under expansion it is minus the
reverse KL divergence. For non-nested supports it is only a signed determinant
or volume contrast and must be accompanied by the enabled/lost/retained ledger.
For a non-counting base measure, the same determinant algebra would instead
define $\chi_x^{(\mu)}$ for normalized restrictions of that measure; it cannot
silently replace the cardinality-defined $\Delta I_x$. The question mark
in Eq.~\eqref{eq:gfe-iota-operator-bridge} carries the missing physics:
the operators, common domain, metric meaning, covariance, functional variation,
and continuum regularization remain to be derived from IOTA and GfE.

\paragraph{Conditional scaling and discriminator.}
If a mode-resolved IOTA contribution $\chi_{x,k}$ and a physical mode dictionary
were independently derived, the conditional sectorwise crosswalk would be
\begin{equation}
 \begin{aligned}
 \chi_{x,k}&\widehat{=}c_k:=\delta s_k/\delta v,\\
 \mathcal G_k&=e^{\chi_{x,k}},\qquad
 \theta_k=e^{\chi_{x,k}}-1,\qquad
 \delta\epsilon_k/\delta v=e^{\chi_{x,k}}-1-\chi_{x,k}.
 \end{aligned}
 \label{eq:gfe-conditional-crosswalk}
\end{equation}
Here $k\in\{0,\ldots,4\}$ labels the five GfE sectors and $z_k$ is the
corresponding degeneracy, with $z_0=z_1=1$ and $z_2=z_3=z_4=3$. In the
reported weak-curvature Friedmann expansion, $u_k=\omega_kH^2$ with
$\omega_k=\omega'_k\ell_{\mathrm P}^2$, and
$\bar\omega_{[1]}:=\sum_{k=0}^{4}z_k\omega_k$.
The signed $\chi_{x,k}$ makes positive sectors conditionally interpretable as
pruning and negative sectors as expansion. This is not a general KL
identification: non-nested mode supports can mix both, and an aggregate
$\chi_x\widehat{=}\sum_kz_kc_k$ determines neither a single temperature nor
$\delta\epsilon/\delta v=\sum_kz_k(e^{c_k}-1-c_k)$ without the full mode vector.

The adjacent-possible bridge also suggests a sharper multiscale signature. A
retained carrier could restrict immediate microstate options while enabling
more long-horizon role types,
\begin{equation}
 \chi_{\mathrm{micro},1}>0,
 \qquad
 \chi_{\mathrm{role},h}<0\quad(h\ge h_\star),
 \label{eq:constraint-enablement-crossover}
\end{equation}
under the same randomized carrier intervention. This local-pruning to
long-range-enablement crossover would become a distinct IOTA/ICP result only if
$h_\star$ were predicted from independently measured gate lifetimes or adaptive
cut depths and confirmed on held-out laws and geometries. Adjacent-possible and
scaffolding theories already anticipate the qualitative phenomenon.

For weak-curvature, late-time Friedmann backgrounds, modes with
$\omega_k\neq0$ obey $c_k\simeq\omega_kH^2$ and
$\delta\epsilon_k/\delta v\simeq\tfrac12\omega_k^2H^4$. For non-de Sitter
solutions with $H\sim t^{-1}$ these scale as $t^{-2}$ and $t^{-4}$. The summed
entropy's leading coefficient $\bar\omega_{[1]}$ vanishes in the reported Milne
and $w=1$ cases, so an aggregate $H^2$ scaling requires
$\bar\omega_{[1]}\neq0$. All of these asymptotics insert standard Friedmann
behavior into the GfE functionals; they are not yet a novel ICP prediction or a
discriminator from GR/\(\Lambda\)CDM.

There is a further conditional scaling crosswalk, but not yet a cosmological
law. If a declared role repertoire grows as $M(v)\propto v^\beta$ and the GfE
mode variable is identified with a constraint potential
$s_C(v):=-\ln M(v)$, then
\begin{equation}
 c(v)=\frac{ds_C}{dv}\simeq-\frac{\beta}{v},
 \qquad
 e^{c(v)}-1-c(v)\simeq\frac{\beta^2}{2v^2}.
 \label{eq:adjacent-gfe-scaling}
\end{equation}
The sign convention matches support expansion in
Eq.~\eqref{eq:signed-support-deformation}. Triggering models already generate
Heaps exponents $\beta$, and TAP models already generate sharp transitions;
Eq.~\eqref{eq:adjacent-gfe-scaling} becomes new physics only after $v$, the mode
dictionary, dimensions, and the carrier-derived rate are fixed independently.

A first bridge-specific test is instead to construct $Q_x$ from frozen IOTA
microdynamics and $\boldsymbol\Theta_x$ independently from the geometric sector,
then test whether their scale-dependent log determinants agree without fitting
the equality. Generic branching, percolation, ordinary complete-state dynamics,
GR/\(\Lambda\)CDM, and unbridged GfE do not themselves supply that cross-level
identity. A cosmological claim requires more: the derived continuum map must
produce a transition, correction, or observable that is absent from the named
complete rivals and is registered before data comparison. At present this is a
precise research bridge, not empirical support for ICP or GfE.

\subsection{Biology supplies persistent records, not a substrate ontology}

\paragraph{Established field result.}
Biological constraint-closure theory models relatively conserved structures
that act on material processes while depending on one another for maintenance
in an open system \citep{MontevilMossio2015Closure}.  The description is
scale- and timescale-relative: an enzyme may constrain one reaction while being
produced and degraded on another horizon.  Whether closure alone grounds
function, normativity, or life remains contested; its most defensible role here
is a model-guided empirical abstraction rather than a definition of purpose
\citep{Garson2017AgainstOrganizational,Corti2023OrganizationalNormativity,
Rider2025ClosureModel}.

Two experiments closely instantiate the persistence question.  In bacterial
acquired immunity, phage exposure installs a sequence record in a CRISPR locus,
and adding or deleting matching spacers changes later resistance
\citep{BarrangouEtAl2007CRISPR}.  In yeast, prior inositol starvation creates
INO1 transcriptional memory involving a DNA element, nuclear-pore association,
histone state, and writer--reader machinery.  The H3K4me2 mark persists for
roughly four divisions after removal of the initiating factor Sfl1, but is lost
when maintenance machinery is removed
\citep{LightEtAl2010INO1,SumpEtAl2022INO1Memory}.  The record outlives its initiator;
it never becomes independent of all material realization.

Evolutionary experiments show a related path effect.  Exhaustive construction
of all 32 combinations of five TEM beta-lactamase substitutions found only 18
of 120 monotonic direct paths accessible under the study's measured landscape
and strong-selection/weak-mutation assumptions
\citep{WeinreichEtAl2006MutationalPaths}.  Ancestral reconstruction of the
glucocorticoid receptor showed that rare permissive substitutions made later
functional changes tolerable \citep{HarmsThornton2014Contingency}.  These are
strong demonstrations of context-dependent probability or operational
accessibility, not context-free proofs of physical impossibility; population
size, drift, recombination, environment, and horizon alter the path graph.

Constructive origin-of-life work supplies both positive and negative controls.
Designed RNA enzymes can cross-catalyze exponential replication, and cooperative
RNA networks can outgrow selfish cycles
\citep{LincolnJoyce2009RNAReplication,VaidyaEtAl2012CooperativeRNA}.  Yet tested compositional
autocatalytic networks show that self-maintenance and recurrence need not yield
selected heredity or evolvability \citep{VasasEtAl2010Evolvability}.  Lipid or
droplet compartments causally alter co-inheritance, resource sharing, and
parasite competition; synthetic protocells can preserve full-length replicators
relative to bulk conditions \citep{AbilEtAl2024ProtocellEvolution}.  Recent synthetic
cells integrate genome replication and lipid synthesis in partial
architectures, while still relying on supplied expression systems, substrates,
energy, and externally imposed cycling \citep{RestrepoSierraEtAl2026SyntheticCell}.

\paragraph{Strongest rival explanation.}
Every example has a sufficient ordinary mechanism: nucleic-acid sequence and
base pairing, chromatin chemistry, protein structure and population dynamics,
reaction kinetics, diffusion, sequestration, and external laboratory
operations.  A state description that omits the spacer, mark, enzyme, boundary,
or bath will appear history-dependent; an adequate physical state includes the
carrier.  Moreover, counterfactual semantic information already scrambles
system--environment correlations and measures effects on later viability
\citep{KolchinskyWolpert2018Semantic}, while Causal Leverage Density proposes
the corresponding comparison over distributions of future phase-space
trajectories \citep{SowinskiEtAl2025CausalLeverage}.  Generic future-shaping is
not ICP novelty.

\paragraph{Synthesis inference and ICP hypothesis.}
Biology strongly supports the architecture
\begin{equation}
 \text{history}\rightarrow\text{embodied record}
 \rightarrow\text{changed later response},
\end{equation}
including cases in which the initiating event or molecule is gone.  The
defensible ICP extension is to add formation, sufficiency, carrier
intervention, succession, and representation-robustness tests across a family
of challenges.  Biology does not establish that information is ontologically
prior to the molecules, nor that IOTA produces living organization.  It does
provide well-controlled physical systems in which the full carrier criterion
can be tested prospectively.

\subsection{Black-hole information is a frontier test of access}

Black-hole theory is relevant because it separates claims often compressed
into the sentence ``the information comes out.''  The necessary hierarchy is
\begin{equation}
 \begin{split}
 \text{global preservation}
 &\not\Rightarrow \text{algebra-relative distinguishability}
 \not\Rightarrow \text{logical reconstruction}\\
 &\not\Rightarrow \text{message recovery}
 \not\Rightarrow \text{efficient physical recovery}.
 \end{split}
 \label{eq:black-hole-hierarchy}
\end{equation}
Each arrow requires a subsystem or algebra, operation class, code subspace,
error tolerance, and resource budget.

\paragraph{Established or controlled field results.}
Hawking's semiclassical calculation motivates pure-to-mixed evolution after an
inaccessible region is traced out \citep{Hawking1976Breakdown}.  Page's curve
is a conditional random-state benchmark: if formation and evaporation are
unitary and the bipartite state is typical, the entropy of a chosen radiation
subsystem rises and falls \citep{Page1993Information}.  Hayden--Preskill
recovery shows rapid message availability under idealized unitary mixing and
access to early radiation, while the Harlow--Hayden result separates
information-theoretic recoverability from feasible decoding by an exponential
complexity barrier \citep{HaydenPreskill2007Mirrors,
HarlowHayden2013Computation}.

Quantum extremal surfaces and islands extremize generalized entropy, combining
an area term with bulk entropy.  In controlled AdS/CFT, JT-gravity-plus-bath,
braneworld, and related models, saddle transitions reproduce Page behavior and
move a logical interior algebra into a radiation entanglement wedge
\citep{EngelhardtWall2015QES,Penington2020EntanglementWedge,
AlmheiriEtAl2019BulkEntropy}.  Holographic quantum error correction and
bulk--boundary relative-entropy results make redundant reconstruction precise
inside a declared code subspace
\citep{AlmheiriDongHarlow2015QEC,JafferisEtAl2016RelativeEntropy}.  Replica
wormholes supply controlled entropy calculations, with ensemble interpretation
and analytic-continuation assumptions kept explicit
\citep{PeningtonEtAl2022ReplicaWormholes,
AlmheiriEtAl2020ReplicaWormholes}.

These are major theoretical advances.  They do not yet provide a unique
microscopic carrier, an efficient decoder for a realistic four-dimensional
evaporating black hole, or direct astronomical observation of a Page curve or
information-bearing Hawking correlations.  A replica saddle is a contribution
to a calculation, not by itself a real-time memory trajectory or a physical
counterfactual branch.

\paragraph{Observable-algebra boundary and current frontier.}
Gravity resists a naive tensor factorization because gauge constraints,
gravitational dressing, and asymptotic charges link nominal regions.  One line
of work argues that a complete asymptotic algebra can distinguish the state and
that the ordinary split property fails \citep{Raju2022SplitProperty}; another
constructs proper radiation subalgebras and Page behavior in massless gravity
under perturbative assumptions \citep{AntoniniEtAl2025Apologia}.  Recent
responses argue that excluding the Hamiltonian or using a blind detector makes
the Page curve a statement about restricted access rather than late emergence
of information \citep{GengEtAl2026Blinders}.  Complementary 2026 models obtain
late saddles or asymptotic recovery in specifically declared massless-gravity
or two-dimensional settings, without closing the full realistic evaporation
problem \citep{Kumar2026MasslessIslands,DarEtAl2026CGHSAlgebras}.  The dispute
cannot be settled by calling both sides ``constraint''; they choose different
observable algebras and therefore answer different operational questions.

Candidate carriers are likewise conditional.  High-order radiation
correlations, code embeddings, gravitational dressing, soft charges, remnants,
and state-dependent interior maps are distinct proposals.  Soft charges are
real conserved structures, but infrared factorization arguments challenge
their capacity to encode hard-state information
\citep{HawkingPerryStrominger2016SoftHair,
MirbabayiPorrati2016Shaving}.  Recent island-encoding work gives a concrete
Goldstone/Stueckelberg dressing mechanism for massive-gravity bath models, not
yet for an isolated astrophysical black hole \citep{Geng2026IslandEncoding}.

\paragraph{Synthesis inference and ICP hypothesis.}
The black-hole lesson is operational modesty.  The same Page-shaped curve can
coexist with different channels, algebras, reconstruction maps, carriers, and
recovery costs.  ICP usefully insists that preservation and accessibility be
indexed by these choices.  It does not currently solve the information paradox
or derive gravity.  A genuinely relevant carrier test would vary a bath
coupling, radiation erasure, detector algebra, code embedding, dressing, or
decoder budget and separately score entropy, distinguishability,
reconstruction error, message fidelity, and cost in one fixed microscopic
theory.

\subsection{Cross-domain result}

The four stress tests converge on a negative and a positive conclusion.  The
negative conclusion is that a Zipf exponent, an entropy value, a persistent
biological response, or a Page curve cannot identify a unique informational
carrier or mechanism.  Native statistics, thermodynamics, molecular dynamics,
quantum channels, and gravity remain sufficient explanations within their
domains.  The positive conclusion is that each field already requires some
version of an indexed possibility structure: a path law, operation preorder,
observable algebra, carrier-bearing state, or resource-bounded decoder.

That convergence supports ICP as a disciplined comparative research program.
It does not establish ICP as a fundamental law.  The proposal gains physical
content only when an interaction-formed role survives original-carrier
replacement and its intervention produces a robust future-law difference not
exhausted by the complete native model.

\section{Adversarial viability audit}

The synthesis can now answer its first question. The answer is neither that
information language is empty nor that IOTA has been confirmed as fundamental.
The evidence supports a narrower conclusion: IOTA/ICP is a coherent,
operationally testable research architecture with bounded constructive results
and independent precedents for its components. Whether it is a distinct
physical principle remains unresolved.

\subsection{The complete-state challenge}

The strongest null hypothesis is ordinary and severe. Let $Z_t$ be a complete
physical specification at time $t$, including all degrees of freedom relevant
at the chosen precision, their correlations, boundary conditions, accessible
resources, controls, and the applicable dynamics. Then the future family is
already fixed as
\begin{equation}
  \Efamily=\Phi_{\tau}(Z_t;\Pi,\mathcal A,B,C).
  \label{eq:complete-state-null}
\end{equation}
On this account, a memory, scar, environmental record, regulatory state, or
IOTA bridge role matters only because it is part of $Z_t$. History is relevant
exactly when it has left a present physical difference; enlarging the state
until it is sufficient removes any independent history variable. Calling that
difference ``information'' and its effects ``constraint'' adds a useful
description, but no new law.

This null is not defeated by showing that two visibly similar systems behave
differently, by finding mutual information between past and future, or by
demonstrating that an environmental degree of freedom carries a record. Hidden
state, conserved quantities, unmeasured correlations, ordinary free energy,
and control mismatch must be exhausted first. Nor is it defeated by a compact
description alone: a predictive statistic may be epistemically economical
without being a new physical entity.

ICP becomes distinct only if its carrier-equivalence construction does useful
work beyond Eq.~\eqref{eq:complete-state-null}. Candidate achievements include
a minimal physically manipulable sufficient class, an invariant across genuine
carrier succession, a new conversion or impossibility theorem, a large and
validated compression of intervention--response laws, or a preregistered
prediction missed by the best complete-state model. The present project has
not yet delivered one at fundamental-physics scope.

The exact IOTA relative-entropy identity and the proposed GfE operator map make
the cross-domain claim sharper, but formal coherence alone does not change this
verdict. The map must first produce a representation-stable invariant or a
preregistered prediction beyond Eq.~\eqref{eq:complete-state-null}; otherwise it
is a calibrated redescription of two native constructions.

The adjacent-possible review increases this pressure. Kauffman's one-step
novelty frontier, triggered urns, evolving discovery networks, scaffolding, and
TAP already cover much of the language of possibility-space expansion. Their
formal alignment with Eq.~\eqref{eq:adjacent-frontier} is a real structural
correlation, but overlap is not independent evidence. The remaining question is
whether a physically identified successor carrier yields a preregistered
frontier change that those complete native models do not.

\subsection{Nine failure modes that remain live}

\begin{enumerate}
  \item \textbf{State-augmentation collapse.} A claimed historical effect
  disappears when the carrier, environment, or multi-time state is adequately
  represented. This would leave ICP as an organizing vocabulary for ordinary
  memory rather than an additional state principle.

  \item \textbf{Representation dependence.} The carrier or constraint exists
  only under a convenient partition, graph encoding, observable algebra, or
  coarse-graining. Quantum factorization results make this a substantive risk;
  the negative IOTA graph-role gate shows it is also a project-level risk.

  \item \textbf{Inferential--physical equivocation.} An observer's loss of
  uncertainty is mistaken for a change in the system's physical possibilities.
  Shannon information, Bayesian conditioning, and a record readout do not by
  themselves establish causal exclusion or a transported role.

  \item \textbf{Generic-dynamics redescription.} Any present state constrains
  its future under a law. Without the formation, intervention, succession, and
  robustness gates, the proposed principle collapses into this truism or into
  existing reachability, viability, semantic-information, or predictive-state
  theory.

  \item \textbf{Engineered-lattice artifact.} IOTA effects may depend on one
  update rule, scheduler, topology, boundary, or analyst-defined address
  scheme. A successful construction establishes possibility within that model,
  not use by nature or universality across realizations.

  \item \textbf{Imported emergence.} A calculation can reproduce a target
  only because quantum amplitudes, a Born rule, a clock, a target Hamiltonian,
  or a mass schedule were supplied. The current Dirac result openly occupies
  this category: valuable role recovery, but not base-level emergence.

  \item \textbf{Summary-statistic underdetermination.} Zipf exponents, entropy
  values, Page curves, and aggregate correlations can be shared by mechanisms
  with different carriers and future laws. A suggestive scalar cannot identify
  the proposed causal architecture.

  \item \textbf{Novelty-by-label or fixed-space discovery.} A first visit to a
  prestated state, graph node, or analyst-defined type is mistaken for creation
  of a new physical possibility. Without a frozen type map and matched
  enabled/lost/retained ledger, an apparent adjacent frontier can be a
  classifier or exploration artifact.

  \item \textbf{Self-confirming evidence.} Project papers, simulations, and
  analysis pipelines may share assumptions and error modes. The disclosed
  Dirac predicate withdrawal and stopped IOTA confirmation are reasons to
  preserve adversarial audits and external replication, not reasons to erase
  the intact bounded results.
\end{enumerate}

\subsection{What the evidence does and does not jointly license}

Several claims survive when the domains are combined. Interactions can create
records in secondary carriers. Present physical organization can alter support,
probabilities, viability, convertibility, work extraction, and accessible
operation order. Some roles survive turnover of their first local realization.
Collective quantum systems can support excitations with particle-like
properties that are not properties of their microscopic constituents. These
are established field results or bounded project constructions within declared
assumptions.

Adjacent-possible research additionally establishes that novelty can be
history-correlated, that exploration can change later accessibility, and that
network, resource, and scaffolding structure alters discovery. It does not
establish that the relevant physical state space is literally created, that
IOTA supplies its mechanism, or that combinatorial growth is a cosmological
law. A Fisher-volume ``biological cosmological constant'' now also makes
geometric expansion and a phase-transition interpretation explicit prior art;
it does not provide ICP's signed carrier intervention
\citep{FarinaVallejo2026LambdaB}.

The conjunction does not establish the IOTA ontology. No reviewed experiment
shows information existing without physical realization; ICP does not require
that stronger claim. More importantly, no current derivation begins with IOTA's
base events and relations and obtains quantum composition and probabilities,
then interacting fields, mass--energy invariants, geometry, and clocks without
placing those structures in the premises. Black-hole reconstruction, Quantum
Darwinism, feedback work extraction, biological organization, and
quasiparticles each demonstrate only a segment of that ladder.

The name \emph{Information--Constraint Principle} remains defensible if
``principle'' denotes the proposed organizing and research principle now under
test. It should not be presented as an established law. Renaming it merely to
avoid that distinction would not strengthen the science; explicit status labels
do.

\subsection{Graded verdict}

Table~\ref{tab:verdict} states the resulting graded verdict and the evidential
boundary attached to each component of the program.

\begin{table}[ht]
\centering
\small
\caption{Viability verdict at the 20 July 2026 bounded-search cutoff, with the
targeted 21 July GfE and adjacent-possible amendments. Ratings refer to the
research program, not to the truth of the ontology.}
\label{tab:verdict}
\begin{tabularx}{\linewidth}{>{\raggedright\arraybackslash}p{0.19\linewidth}
  >{\raggedright\arraybackslash}p{0.15\linewidth}
  >{\raggedright\arraybackslash}X}
\toprule
Dimension & Rating & Evidential basis and ceiling \\
\midrule
Conceptual coherence & Strong & Retained condition, informational change,
intervention, future family, and carrier succession are separable and
operationally specifiable. \\
Operational precision & Moderate--strong & The four carrier gates and matched
future-family contrasts can fail; the signed flux ledger handles pruning and
expansion, but a generally accepted invariant and physical measure are absent. \\
Distinctness from prior art & Partial & Formation plus succession beyond the
original witness is narrower than generic reachability and adjacent-possible
expansion, but complete-state, causal-state, process-tensor, and
semantic-information accounts remain live. \\
Project evidence & Moderate, bounded & Frozen IOTA interventions establish
constructive examples; Dirac role recovery is intact within its supplied
quantum layer. Neither is an independent observation of nature's substrate. \\
Independent grounding & Indirect, substantial & Physics and biology establish
component mechanisms, not their IOTA integration. \\
Falsifiability & Strong locally & Carrier edits, state augmentation,
representation changes, and held-out interventions can reject individual
claims; no single finite test yet rejects every possible IOTA realization. \\
Information-first ontology & Open & No result uniquely selects informational
events and relations as nature's substrate or derives all supplied physics. \\
Research tractability & Strong & The next bridges have concrete observables,
controls, nulls, and stopping rules. \\
\bottomrule
\end{tabularx}
\end{table}

The resulting verdict is positive but graded:

\begin{quote}
\textbf{IOTA/ICP has real scientific legs as a coherent, falsifiable research
program and as an operational synthesis of persistent constraint. Bounded
project results demonstrate beginnings of the proposed mechanism. Current
science supplies independent precedents for its components, but neither their
integration nor an information-first ontology has been established. ICP is
therefore a viable proposed principle and research architecture, not an
established fundamental law.}
\end{quote}

This verdict is more useful than an unconditional endorsement. It identifies
exactly what survives current science and the observations that would force a
weaker interpretation.

\section{Research path: from persistence to physics}

The synthesis's second function is to convert the verdict into an ordered
program. The order matters. A higher rung cannot be counted as evidence for a
lower rung when the lower structure was supplied in advance.

\subsection{Rung 1: replicate carrier succession}

The immediate goal is not a more elaborate cosmological simulation. It is a
harder test of the result already obtained. Repeat the disjoint-replacement and
lineage protocols under alternative local laws, asynchronous and randomized
schedulers, topologies, boundary conditions, state encodings, and independent
implementations. Freeze the role equivalence and opportunity horizon before
opening outcomes. Use held-out source geometries and publish null and censored
cases. Estimate Eq.~\eqref{eq:icp-constraint-divergence} under the same frozen
choices and require its qualitative result to survive law, scheduler, topology,
encoding, boundary, and coarse-graining changes before treating it as more than
a model-indexed diagnostic.

Success would establish that constraint succession is a class of mechanism
rather than a quirk of one coordinate system. Failure under modest
representation or rule changes would lower the claim to a model-specific
phenomenon. Either outcome is informative.

\paragraph{Adjacent-frontier assay.}
On each held-out system, freeze the type map $r_\ell$, actualized repertoire,
policy family, resources, and horizons before running interventions. Enumerate
or bound $\mathcal J_{t,\tau}^{\pi,\ell}$ and report enabled, lost, and retained
identities rather than novelty counts alone. Randomize carrier preservation,
ablation, restoration, and sham edits; compare a fixed-state Markov model,
reinforced urn, evolving-network, branching/percolation, and TAP rivals. The
strongest near-term signature is not another Heaps or Zipf fit. It is a
prospectively predicted crossover in
Eq.~\eqref{eq:constraint-enablement-crossover}: local pruning followed by
greater long-horizon role-level accessibility, with $h_\star$ forecast from
separately measured opportunity lifetimes or adaptive veto depth.

\subsection{Rung 2: isolate a minimal persistent coherence}

A connectivity role is not yet a quantum coherence. The next target should be
the smallest organization whose empirically complete future family cannot be
compressed to any one constituent or original address, yet is preserved by a
registered succession map. Required measurements include lifetime,
propagation, composition, interference-like response, robustness to local
perturbation, and a causal edit of the current carrier. Competing ordinary
state summaries must be fit on training interventions and evaluated on held-out
interventions and horizons.

A decisive result would identify a representation-robust equivalence class of
realizations with a common intervention--response role. A decisive negative
would show that each apparent coherence reduces to a changing list of local
states with no stable role beyond the analyst's tracking rule.

\subsection{Rung 3: derive the quantum layer}

The quantum bridge must produce, rather than assume, the structures used by the
Dirac construction. At minimum it must account for:
\begin{enumerate}
  \item a state space with a defended composition rule;
  \item reversible or appropriately open effective dynamics;
  \item interference and contextual probability constraints;
  \item a measurement or intervention semantics;
  \item the Born probability rule, or a quantitatively distinct replacement;
  \item operational clocks and reference-frame transformations; and
  \item a controlled continuum or scaling regime.
\end{enumerate}
Merely storing complex numbers on lattice sites would import the target. The
derivation must show why the effective equivalence classes and transformation
rules arise from base succession and why relevant predictions are insensitive
to arbitrary partition or coordinate choices.

The strongest near-term physical analogue is a carrier-edit experiment. Form a
record through two matched histories, characterize its secondary physical
carrier, then preserve, erase, swap, scramble, or rewrite that carrier while
balancing energy, charges, fields, pulses, and environmental correlations.
Estimate the later process or reachable family under a common intervention set.
If an augmented standard quantum state predicts every result, ICP gains a
useful embodied-memory instance but no extra ontology. A preregistered residual
beyond that model would materially raise the verdict.

\subsection{Rung 4: construct a quasiparticle role}

Once a quantum layer exists, the appropriate target is a manipulable collective
excitation, not the full Standard Model. The registered role should require:
creation and annihilation procedures; localized or relational transport;
dispersion; composition; symmetry transformation; a conserved or topological
label; controlled scattering; robustness and lifetime; and convergence across
scales. Established string-net, topological-order, quantum-simulator, and QCA
results prove that such emergence is possible in other local quantum systems.
They also supply hostile benchmarks and lattice no-go constraints.

The current Dirac construction then becomes a downstream validation target: a
derived excitation should reproduce its bounded operational witnesses without
being tuned by an externally supplied Dirac map, mass schedule, or Born
readout.

\subsection{Rung 5: interacting fields, mass, and energy}

Free one-particle kinematics is not a field theory. The next stage must recover
many-excitation composition, statistics, interactions, locality or a precise
replacement, scattering observables, conservation laws, and controlled
renormalization or continuum behavior. Fermion doubling, symmetry recovery,
gauge constraints, anomalies, and Lorentz-violation bounds are stopping tests,
not optional finishing details.

Only after this stage should the project claim emergence of mass and energy.
The target is not an analogy between information and work. It is an effective
invariant that plays the mass or energy role across processes: governing
dispersion, time translation, inertia, conservation and exchange, and coupling
in the appropriate regime. Its numerical value and transformation law must be
derived from the organization, not placed in an update angle or external
Hamiltonian.

\subsection{Rung 6: geometry, clocks, and macroscopic recovery}

An information-first theory must ultimately explain why effective causal and
metric relations approximate the spacetime structure already tested by
physics. The project must identify operational rods and clocks, recover
frame-covariant predictions, specify when causal order is definite, and derive
the scale at which continuum geometry is accurate. Black-hole information is a
late stress test of that theory---especially algebra-relative recovery,
generalized entropy, and computational accessibility---not an early proof of
the substrate.

Macroscopic recovery also requires thermodynamic and statistical explanation.
The theory must say why entropy concepts appropriate to different tasks agree
where they should, diverge where they must, and yield ordinary irreversible
phenomena without turning one entropy scalar into the definition of
constraint. Zipf-like distributions, if they occur, require a mechanism and
matched-marginal counterfactual tests rather than a fitted slope.

The GfE bridge in Eq.~\eqref{eq:gfe-iota-operator-bridge} makes one part of this
rung unusually concrete. Construct $Q_x$ from IOTA future laws rather than from
the target geometry; prove positivity and covariance; control the finite-size
and continuum determinants; and test Eq.~\eqref{eq:gfe-iota-scalar-bridge}
across scales and held-out geometries. The pruning/expansion orientation in
Eq.~\eqref{eq:signed-kl-duality}, the cross-scale sign change in
Eq.~\eqref{eq:constraint-enablement-crossover}, and the conditional scaling in
Eq.~\eqref{eq:adjacent-gfe-scaling} must be predicted from carrier dynamics,
not chosen after seeing a geometric trend. Only then should the program derive a
cosmological correction, transition scale, or observation and compare it with
GR/\(\Lambda\)CDM, native GfE, the Fisher-volume $\Lambda_B$ construction, and
nongravitational branching or percolation nulls. The immediate goal remains the
harder carrier-succession replication in
Rung~1, not a premature cosmological simulation.

\subsection{Decision gates and stopping rules}

The program should continue while each rung yields preregistered constructive
results, useful failures, sharper invariants, or discriminating predictions.
It should be narrowed or radically reformulated if any of the following
persists after technically adequate attempts:
\begin{itemize}
  \item complete-state augmentation repeatedly removes the candidate
  historical variable;
  \item carrier interventions fail under matched physical controls;
  \item the role disappears under reasonable representation, factorization,
  scheduler, or coarse-graining changes;
  \item the quantum bridge requires Hilbert-space, probability, clock, or
  target-dynamics structure to be inserted by hand;
  \item higher-dimensional, interacting, symmetry, or continuum tests fail
  their relevant no-go boundaries; or
  \item no new compression, invariant, prohibition, or physical prediction
  emerges beyond native theories.
\end{itemize}

These are scientific stopping rules, not rhetorical concessions. A result that
reduces ICP to a precise theory of physically embodied memory would still be
valuable, even if the information-first ontology fails.

\section{Conclusion}

The original question survives the audit in a sharpened form. Information can
persist beyond the object that first carried it, but every established example
does so through continuing physical realization in another state, relation, or
organization. IOTA has constructed bounded cases in which a relational
constraint is destroyed locally and re-instantiated through a causally
necessary successor. Established science shows that records can be distributed,
biological constraints recursively maintained, future options formally
restricted, and particle-like roles produced collectively. Those results make
the larger program scientifically plausible; they do not complete it.

The honest present claim is therefore deliberately finite. IOTA/ICP has begun
to demonstrate the mechanism it proposes and has a tractable path on which the
next steps can succeed or fail. The next result that would materially change
the verdict is not another cross-domain resemblance. The adjacent possible is
an exceptionally close formal neighbor, and the signed support identity turns
its expansion/pruning duality into a falsifiable bridge target while narrowing
generic novelty claims as prior art. No carrier-derived operator map or novel
observation has yet been produced. The decisive next result is still a minimal
persistent
coherence, derived from the base IOTA layer, that survives carrier succession
and supports a held-out quantum intervention role without importing quantum
state structure or probability. From there, a quasiparticle is a defensible
next target. Until those bridges are crossed, information primacy remains the
program's hypothesis rather than its conclusion.


\section*{Data and code availability}
The literature registry, exact search log, claim--evidence matrix, rival and
falsification registers, formal derivations, and manuscript sources are
available at \url{https://github.com/fmertz55/ICP}. The versioned simulation
code, preregistrations, evidence manifests, and frozen outputs supporting the
IOTA project results are available at
\url{https://github.com/fmertz55/iota}.

\section*{Ethics statement}
This theoretical and computational study used no human participants, human
tissue, animals, or field samples; ethics approval and consent were not
applicable.

\section*{Funding}
This research received no external funding.

\section*{Declaration of competing interests}
The author declares no known competing financial interests or personal
relationships that could have appeared to influence the work reported here.

\section*{CRediT author statement}
Douglas Clark: Conceptualization; Methodology; Formal analysis; Investigation;
Data curation; Software; Validation; Visualization; Writing -- original draft;
Writing -- review and editing; Project administration.

\section*{Declaration of generative AI and AI-assisted technologies}
During preparation of this work, the author used OpenAI Codex for
AI-assisted literature triage, drafting, code review, and LaTeX editing. The
author reviewed and verified the cited sources, calculations, code, claims, and
final text and takes full responsibility for the content of the publication.

\bibliographystyle{plainnat}
\bibliography{library/references}

\end{document}
